\documentclass[12pt,letterpaper,boxed]{hmcpset}
\usepackage[margin=1in,headheight=14pt]{geometry}
\usepackage{amsfonts, amsmath, amssymb, enumerate, fancyhdr, gensymb, lastpage, mathtools, parskip, graphicx}
\usepackage{xcolor, tikz-cd}
\newcommand{\wg}[1]{\textcolor{violet}{#1}}
\newcommand{\OO}{\mathcal O}
\newcommand{\Q}{\mathbb Q}
\newcommand{\R}{\mathbb R}
\newcommand{\C}{\mathbb C}
\newcommand{\Z}{\mathbb Z}
\renewcommand{\H}{\mathbb H}
\newcommand{\abs}[1]{\left|#1\right|}
\newcommand{\abssq}[1]{\left|#1\right|^2}
\newcommand{\im}{\text{im }}
\newcommand{\inv}{^{-1}}
\newcommand{\normal}{\unlhd} %% one can also use \trianglelelefteq
\newcommand{\anglee}[1]{\langle #1 \rangle}
\usepackage[shortlabels]{enumitem}

% Numbering macros
\pagestyle{fancy}
\lhead{Will Gilroy}
\chead{Complex Homework \#12}
\rhead{06 May 2026}
\lfoot{}
\cfoot{}
\rfoot{Page\ \thepage\ of\ \pageref{LastPage}}

\linespread{1.5}

\newcommand\blankpage{
    \thispagestyle{empty}
    \addtocounter{page}{-1}
    \newpage}
\renewcommand\footrulewidth{0.4pt}

\begin{document}

\problemlist{Complex Analysis Homework \#12} 

%------------------------- Problem 1 -----------------------

\begin{problem}
	\includegraphics[scale=0.8]{1.png}
	\hfill
\end{problem}

\begin{solution}
\begin{enumerate}
\item Let us write $\omega_1/ \omega_2 = p/q$ for some coprime
integers $p,q$. 
Let us define $\omega_0 := \omega_1/p$.
We need to show that $(1)$ $f$ is periodic with repsect to $\omega_0$
and $(2)$ that $\omega_1, \omega_2$ are multiples of $\omega_0$. 

We show $(2)$ first. Notice that $\omega_1 = p\omega_0$ by definition.
And also, $\omega_2 = \frac{q\omega_1}{p} = q \omega_0$. 

Next we show $(1)$, that $f$ is periodic with repsect to $\omega_0$. 
Recall B\'ezout's identity gives us that there exists integers $m,n$
such that $mp + nq = 1$ since $p,q$ are coprime. 
Rearragnging gives us that 
\begin{align*}
	\frac{\omega_1}{p}
		&= \omega_1\left(m + n \frac{q}{p}\right) \\
		&= \omega_1\left(m + n \frac{\omega_2}{\omega_1}\right) \\
	\omega_0 &= m\omega_1 + n \omega_2.
\end{align*}
Since $f$ is periodic in $\omega_1$ and $\omega_2$ it thus follows
that $f$ is periodic in $\omega_0$. 

\item Let $\tau := \omega_1/\omega_2$, we have that $\tau$ is
irrational.
Let $m,n \in \Z$ and notice that $m\omega_1 + n \omega_2 =
\omega_2(m\tau + n)$. Now since $\tau$ is irrational we have that the
set $\{\omega_2(m \tau + n)\}$ is dense in the line $\R\omega_2$. 
We will show that $f$ is thus constant on the line $\R\omega_2$. 

Let $c \in \R\omega_2$ be a point where $f$ is holomorphic.
Since $\{\omega_2(m\tau + n)\}$ is dense in $\R \omega_2$
we have a sequence $z_k = \omega_2(m_k \tau + n_k)$ with $m_k,n_k \in \Z$ which converges to $c$. 
Since $f$ is holomorphic at $c$ we have that it is in particular
continuous at $c$, and so $f(c) = \lim_{k \to \infty}f(z_k)$. 
Moreover since $\omega_2(m_k\tau +n_k) = m_k\omega_1 + n_k\omega_2$
for all $k$ we have
that $f(\omega_2(m\tau+n)) = f(0)$ for all $m,n \in \Z$. 
But then $f(c) = \lim_{k \to \infty} f(z_k) = f(0)$, i.e. $f$ is constant along
$\R\omega_2$.

We have shown that $f$ is constant on a closed subset of its domain of holomorphicity.
It follows by the identity theorem that $f$ is thus constant on its
entire domain.




\end{enumerate}
\end{solution}

\newpage

%------------------------- Problem 2 -----------------------

\begin{problem}
	\includegraphics[scale=0.8]{2.png}
	\hfill
\end{problem}

\begin{solution}
Since $F$ is a meromorphic function on $\C/\Lambda$ it lifts to a
doubly periodic meromorphic function on $\C$, let us also refer to
this meromorphic function on $\C$ as $F$ (I will note when we start
thinking about it as a function on the torus again). 
Consider a fundamental paralellogram $P$ of $\Lambda$ in $\C$ such that
$F$ has no zeroes or poles on the boundary of $P$. 
Recall that the following variation of the argument principle integral
counts the sum of the points which gives zeroes and poles for $F$.
\[
	\int_{\partial P} z \cdot \frac{f'(z)}{f(z)} dz = \sum_{z_i \in \C} z_i - \sum_{p_j \in \C} p_j.
\]

Let us normalize our lattice so that one of the generators is given by
$1$ and the other is given by $\tau \in \mathbb H$. 
Let us label the sides of our fundamental paralellogram by $\alpha,
\alpha', \beta, \beta'$.
Then notice that we have the following identity 
\begin{align*}
	\int_{\beta} z \cdot \frac{f'(z)}{f(z)} dz 
	&= \int_{-\beta' + \tau} z \cdot \frac{f'(z)}{f(z)} dz \\
	&= \int_{-\beta'} (z + \tau) \cdot \frac{f'(z+\tau)}{f(z+\tau)} dz \\
	&= \int_{-\beta'} (z + \tau) \cdot \frac{f'(z)}{f(z)} dz \\
	&= -\int_{\beta'} (z + \tau) \cdot \frac{f'(z)}{f(z)} dz,
\end{align*}
where the third equality follows since $f$ is periodic in $\tau$ and
so $f'$ is also periodic in $\tau$. 
Similar reasoning gives that 
\[
	\int_{\alpha} z \cdot \frac{f'(z)}{f(z)} dz = -\int_{\alpha'} (z+1) \cdot \frac{f'(z)}{f(z)}.
\]
Now, consider the following
\begin{align*}
	\int_{\partial P} z \cdot \frac{f'(z)}{f(z)} dz 
	&= \int_{\alpha + \beta + \alpha' + \beta'} z \cdot \frac{f'(z)}{f(z)} dz \\
	&= \int_{\alpha'} \frac{f'(z)}{f(z)} dz - \tau \int_{\beta'} \frac{f'(z)}{f(z)} dz.
\end{align*}
Now, we claim (as prompted by the book hint) that the integral $\int
f'(z) / f(z) dz$ along any of the side lengths of the fundamental
parallelogram is some multiple of $2\pi i$. 
Consider 
\[
	\int_{\alpha'} \frac{f'(z)}{f(z)} dz 
	= \int_{z_0}^{z_0 + \tau} \frac{f'(z)}{f(z)} dz 
	= \log(f(z_0)) - \log(f(z_0 + \tau)),
\]
where $z_0$ is one endpoint for $\alpha'$. Note that, 
since we are integrating just over an interval in $\C$, we can
take the branch cut of $\log$ to be away from either of the line segments
that we care about. 
Now recall that $f$ is periodic with respect to $\tau$, i.e., $f(z_0)
= f(z_0 + \tau)$. 
That means, 
the image of $f$ under $\alpha'$ traces out some closed curve in
$\C^\times$ (recall that we chose $P$ so that $f$ does not vanish on
$\partial P$). 
It follows then that 
\[
	\int \frac{f'(z)}{f(z)} dz = k(2\pi i),
\]
for some integer $k$. 

Now returning the the integral that we are concerned with, we thus have 
\[
	\int_{\alpha'} \frac{f'(z)}{f(z)} dz - \tau \int_{\beta'} \frac{f'(z)}{f(z)} dz
	= (2\pi i)(k_1 - k_2 \tau).
\]
For some $k_1, k_2 \in \Z$. 
Now notice that $k_1 - k_2 \tau$ is a lattice point. And so 
$k_1 - k_2 \tau \equiv 0 \bmod \Lambda$. 
In other words, we have shown 
\[
	\sum z_i - \sum p_j = 
	\int_{\partial P} z \cdot \frac{f'(z)}{f(z)} dz \equiv 0 \bmod \Lambda.
\]


\end{solution}

\newpage

%------------------------- Problem 3 -----------------------

\begin{problem}
	\includegraphics[scale=0.8]{3.png}
	\hfill
\end{problem}
\begin{solution}
Recall that without loss of generality we can take our lattice $\Lambda$ to be
generated by $1$ and $\tau = a + bi$ with $a,b \in \R$ and $b > 0$. We can also choose
generators of our lattice so that $a > 0$. 
With this notation we have that $\omega = m + n \tau$ for integers
$m,n$, and so 
\[
\abssq{\omega} = \abssq{m + n(a+bi)} = (m+na)^2 +(nb)^2.
\]
Our series of interest becomes 
\[
	\sum_{\omega \in \Lambda^\star} \frac{1}{\abssq{\omega}}
	= 
	\sum_{m,n \in \Z } \frac{1}{(m+na)^2 + (nb)^2}.
\]
For $m,n$ not both simultaneously zero. Consider this sum for some
fixed $n$. We can split our sum 
\begin{align*}
	\sum_{m,n \in \Z } \frac{1}{(m+na)^2 + (nb)^2} &= 
	\sum_{(m+na)^2 \leq (nb)^2 } \frac{1}{(m+na)^2 + (nb)^2}  
	+ \sum_{(m+na)^2 > (nb)^2 } \frac{1}{(m+na)^2 + (nb)^2} \\
	& \geq 
	\sum_{(m+na)^2 \leq (nb)^2 } \frac{1}{(m+na)^2 + (nb)^2} \\
	& \geq \sum_{(m+na)^2 \leq (nb)^2} \frac{1}{2(nb)^2} \\
	&= K(n) \cdot \frac{1}{2(nb)^2}.
\end{align*}
Where $K(n)$ is the number of integers $m$ satisfying $(m+na)^2 \leq (nb)^2$.
The second line follows since the latter sum on the first line
contains only 
positive terms. The third line follows since, by the condition in the
sum, $(nb)^2 \geq (m+na)^2$ for each term.
Counting the number of integers $m$ satisfying $(m+na)^2 \leq (nb)^2$
gives us that $K(n) \sim 2\abs{nb}$ and so we have 
\[
	\sum_{m,n \in \Z } \frac{1}{(m+na)^2 + (nb)^2} \geq \frac{1}{\abs{nb}}.
\]
Now we can show that our series of interest diverges 
\[
	\sum_{\omega \in \lambda} \frac{1}{\abssq{\omega}} = 
	\sum_{m,n \in \Z} \frac{1}{(m+na)^2 + (nb)^2}
	\geq \sum_{n \in \Z} \frac{1}{\abs{nb}} = \frac{1}{b} \sum_{n \in \Z} \frac{1}{\abs n},
\]
which diverges.

\end{solution}

\newpage

%------------------------- Problem 4 -----------------------

\begin{problem}
	\includegraphics[scale=0.8]{4.png}
	\hfill
\end{problem}

\begin{solution}
We approach this problem by attempting to construct the desired
Riemann map and its inverse. 


First we construct a candidate Riemann map $\H \to S$ with the
prescribed images.
As a first approach,
let $f: \H \to S$ via $f(z) = -i\log(iz)$, with the branch cut in the
lower right quadrant of $\C$.

Consider $f: \H \to \C$ via $f(z) = -i\log(iz)$
with the branch cut in the lower half plane.
Notice $f(1) = -i\log(i) = -i\log(e^{i\pi/2}) = -i(i\pi/2) =
\pi/2$. Likewise, a similar calculation will give that $f(-1) =
-\pi/2$ and $f(\infty) = \infty$. 
Since we chose the branch cut appropriately, we have that $f$ is
holomorphic on $\H$. Consider the image of $f$. 
Let $z = re^{i\theta} \in \H$ so that $r > 0$ and $0 < \theta < \pi$. 
Now
\begin{align*}
	f(z) &= -i \log(i re^{i\theta}) \\ 
	&= -i \log(r e^{i(\theta + \pi/2)}) \\
	&= -i(\ln(r) + i(\theta + \pi/2)) \\ 
	&=  (\theta + \pi/2) - i \ln(r).
\end{align*}
We have that the image has real part contained in $-\pi/2$ and
$\pi/2$, but the imaginary part takes any value. And so the image of
this map is the non-truncated strip bounded just by the lines $z =
-\pi/2$ and $z = \pi/2$.
Although, $f$ has inverse given by $w \mapsto e^w$.

This calculation may give us an insight the following map, as a second
guess. 
Consider $g(z) : S \to \H$ given by $g(z) = \sin(z)$.
This is holomorphic on $S$ and extends continuously to the boundary since it is
entire as a function $\C \to \C$. 
We need to verify that $g$ is a biholomorphism to $\H$ and that it
takes the prescribed triple to the prescribed image. 

First, we need to verify that the image of $g$ is $\H$. 
Write $z = a + bi$ with $-\pi/2 < a < \pi/2$ and $b > 0$. 
Now consider the following
\begin{align*}
	\sin(a + bi) &= \sin(a)\cos(bi) + \cos(a)\sin(bi).
\end{align*}
Recall the identities
\[
	\sin(x) = \frac{e^{ix} - e^{-ix}}{2i} \qquad \cos(x) = \frac{e^{ix} + e^{-ix}}{2} 
\]
In particular 
\[
	\sin(ib) = \frac{e^{-b} - e^{b}}{2i} \qquad \cos(ib) = \frac{e^{-b} + e^{b}}{2} \in R.
\]
And so we have that the real part of $\sin(a+bi)$ is given entirely by
the second term, 
$Im(\sin(a+bi)) = Im(\cos(a)\sin(bi))$. 
Consider 
\begin{align*}
	\cos(a)\sin(bi) 
	&= \cos(a) \frac{e^{-b} - e^{b}}{2i} \\
	&= \cos(a) (-2i) \frac{e^{-b} - e^{b}}{4} > 0.
\end{align*}
Where the final inequality follows
since $-\pi/2 < a < \pi/2$ gives $\cos(a) > 0$, and $b >
0$ gives $e^{-b} - e^{b} < 0$.
That is \[
	Im(\sin(a+ib)) > 0
\]
And so the image of $g$ is contained in $\H$.

Consider the prescribed triples.
We have $\sin(-\pi/2) = -1$ and $\sin(\pi/2) = 1$. 
Note that $\sin$ does not have a well-defined extension to infinity,
when considered as a function $\C \to \C$. In particular, if we
restrict our attention to the real line, $\sin(x)$ does not grow
unboundedly as $x \to \infty$ along $\R$ (the more formal way to argue
this would be to study what happens with $\lim_{z \to 0}\sin(1/z)$). 
However, since we are considering $\sin$ as a function on $S$, the
only way for $z$ to approach $\infty$ is by sequences $z_n$ where
$-\pi/2 < Re(z_n) < \pi/2$ and $Im(z_n) > 0$ for all $n$. 
In other words, we can only construct sequences that approach infinity
by following the positive imaginary direction. For all such sequences
$\abs{\sin(z_n)}$ does indeed grow without bound (in particular $(e^{-b} -
e^{b})/4 \to -\infty$ as $b \to \infty$). 
Thus we have that $\sin(\infty) = \infty$. 
We have shown that $g$ takes $(-\pi/2, \pi/2, \infty) \mapsto
(-1,1,\infty)$. 

Lastly, we argue that $g$ is a bijection. Notice that the quantity 
$(e^{-b} - e^{b})/4$ is surjective to $\R_-$ for $b>0$, and so $Im(g)$ takes every
positive imaginary value. Likewise $(e^{-b} + e^b) / 2$ is surjective
to $\R_+$ for $b > 0$, moreover, $\sin(a)$ takes each value between
$-1, 1$ as $-\pi/2 < a < \pi/2$. And so $Re(g)$ takes every real
value. We can thus achieve any value in the upper half plane via $g$
by tuning $a,b$. 

We argue that $g$ is injective. Suppose we had $z_1, z_2 \in S$ such
that $g(z_1) = g(z_2)$. Recalling the exponential identity for $\sin$
this becomes the condition 
\[
	k - \frac{1}{k} = m - \frac{1}{m},
\]
where $k = e^{iz_1}$ and $m = e^{iz_2}$. 
This can be rearranged as follows
\begin{align*}
	(k-m) - \left(\frac{1}{k} - \frac{1}{m}\right) &= 0\\
	(k-m)  - \left(\frac{m - k}{mk}\right) &= 0 \\
	(k-m)\left(1 + \frac{1}{mk}\right) &= 0.
\end{align*}
Now if $k = m$ then we have $e^{iz_1} = e^{iz_2}$ which implies 
$iz_1 = iz_2 + 2\pi i \ell$ for some integer $\ell$. However 
this then gives $z_1 = z_2 + 2\pi k$, and no two elements of $S$ can
differ by a nonzero integer multiple of $2\pi i$ since the strip has
real width $\pi$. In other words $k=m$ implies $z_1 = z_2$.
Now, on the other hand, consider the condition $1 + 1/mk = 0$.
We can rearrange this to acquire $mk = -1$. 
First, note that $m,k$ are contained in the right half plane. 
Indeed if $a + bi$ satisfies $-\pi/2 < a < \pi/2$ then
$e^{i (a+bi)} = e^{ia - b} = e^{-b} e^{ia}$ has argument $-\pi/2 < a <
\pi/2$. There exist no $m,k$ in the right half plane which satisfy $mk
= -1$; there are no two angles in $[-\pi/2, \pi/2]$ that we can sum
to $\pi$.  And so the condition $1 + 1/mk = 0$ gives a contradiction.
We have shown that $g$ is injective as a map $S \to \H$. 
We have overall shown that $g$ is a biholomorphism $S \to \H$. 

We have shown that $g(z) = \sin(z)$ is the inverse for the claimed
Riemann map. And indeed this map can be extended to a holomorphic map
on all of $\C$ since $\sin(z)$ is entire. 






\end{solution}

\newpage

\end{document}
